%% Le tableau de conversion des surfaces : deux colonnes par unité, un exemple placé.
%% L'exemple placé (0,29 dm2 = 29 cm2) donne la réponse : il passe par \rappel.
\documentclass[tikz,border=4pt]{standalone}
\usepackage{liaisons}
\begin{document}
\begin{tikzpicture}[x=1cm,y=1cm]
\def\w{1.16}          % largeur d'une unité (deux cases)
\def\h{0.58}          % hauteur d'une ligne

%% en-têtes
\foreach \i/\u in {0/{km$^2$}, 1/{hm$^2$}, 2/{dam$^2$}, 3/{m$^2$}, 4/{dm$^2$}, 5/{cm$^2$}, 6/{mm$^2$}} {
  \pgfmathsetmacro\x{\w*\i+\w/2}
  \node[font=\small\bfseries, text=solideB] at (\x,0.42) {\u};
}

%% grille : trait épais entre unités, trait fin au milieu
\foreach \i in {0,...,6} {
  \pgfmathsetmacro\xa{\w*\i}
  \pgfmathsetmacro\xm{\w*\i+\w/2}
  \draw[line width=0.9pt, draw=black!70] (\xa,0.15) -- (\xa,-2*\h+0.15);
  \draw[line width=0.35pt, draw=black!35] (\xm,0.15) -- (\xm,-2*\h+0.15);
}
\draw[line width=0.9pt, draw=black!70] (7*\w,0.15) -- (7*\w,-2*\h+0.15);
\foreach \j in {0,1,2} {
  \pgfmathsetmacro\y{0.15-\h*\j}
  \draw[line width=0.9pt, draw=black!70] (0,\y) -- (7*\w,\y);
}

%% l'exemple : 0,29 dm2 posé dans le tableau, lu 29 cm2
\rappel{%
  \node[font=\small\bfseries, text=solideA] at (4*\w+0.75*\w,-0.14) {0};
  \node[font=\small\bfseries, text=solideA] at (5*\w+0.25*\w,-0.14) {2};
  \node[font=\small\bfseries, text=solideA] at (5*\w+0.75*\w,-0.14) {9};
  \node[font=\small\bfseries, text=solideA, inner sep=0pt] at (5*\w,-0.28) {,};
  \node[font=\small, text=solideA, anchor=west, inner sep=3pt] at (7*\w+0.1,-0.14)
    {$0{,}29\ \mathrm{dm^2} = 29\ \mathrm{cm^2}$};
}
\node[font=\scriptsize\itshape, text=black!60, anchor=north west, inner sep=2pt]
  at (0,-2*\h+0.05) {deux chiffres par unité : on avance de deux rangs à chaque changement d'unité};
\end{tikzpicture}
\end{document}
